\documentclass[aspectratio=169,11pt]{beamer}
\usetheme{default}
\setbeamertemplate{navigation symbols}{}
\usepackage{graphicx}
\usepackage[T1]{fontenc}
\definecolor{ncsured}{HTML}{CC0000}
\definecolor{ncsudk}{HTML}{990000}
\setbeamercolor{structure}{fg=ncsured}
\setbeamercolor{frametitle}{fg=white,bg=ncsured}
\setbeamercolor{title}{fg=ncsudk}
\setbeamercolor{itemize item}{fg=ncsured}
\setbeamercolor{itemize subitem}{fg=ncsured}
\setbeamerfont{frametitle}{series=\bfseries}
\setbeamertemplate{frametitle}{%
  \nointerlineskip\vskip2pt%
  \begin{beamercolorbox}[wd=\paperwidth,ht=2.4ex,dp=1ex,leftskip=0.3cm]{frametitle}%
  \bfseries\insertframetitle%
  \end{beamercolorbox}}
\setbeamertemplate{footline}{%
  \hbox{\begin{beamercolorbox}[wd=\paperwidth,ht=2.2ex,dp=1ex,leftskip=.3cm,rightskip=.3cm]{}%
  \tiny Optimal Market Making for Cryptocurrency Options \hfill Week 1 \hfill \insertframenumber%
  \end{beamercolorbox}}}

\title{Week 1: Introduction}
\subtitle{Program overview, tooling, and data}
\date{North Carolina State University \,\textbar\, Summer 2026}
\titlegraphic{\includegraphics[width=4.2cm]{/Users/dendisuhubdy/Github/ncsu/assets/ncsu_logo.png}}
\begin{document}
\begin{frame}[plain]\titlepage
\begin{center}\tiny Industry Mentor: Dendi Suhubdy (Bitwyre)  ·  Guest: Fenni Kang (Coincall)\end{center}\end{frame}
\begin{frame}{Welcome to the Program}
\begin{itemize}
  \item A 10-week summer research program on optimal market making for crypto options
  \item Goal: design, implement, and evaluate quoting strategies on real Coincall data
  \item You will produce a research-grade report and a working code repository
  \item This week: orientation, environment setup, and data access
\end{itemize}\end{frame}
\begin{frame}{Maker vs. Taker}
\begin{itemize}
  \item A market maker continuously quotes both sides and earns the spread
  \item A taker crosses the spread to execute immediately
  \item We build the maker's side; in Week 9 a Coincall practitioner covers taker strategies
  \item The spread is compensation for committing to quotes under uncertainty
\end{itemize}\end{frame}
\begin{frame}{The Spine: Avellaneda-Stoikov}
\begin{itemize}
  \item The whole course is organized around the Avellaneda-Stoikov framework
  \item Introduced Week 3; turned into a calibrated engine Week 5
  \item Extended to multi-Greek options inventory Week 8
  \item Advanced extensions (multi-asset, adverse selection, robust) in Week 10
\end{itemize}\end{frame}
\begin{frame}{Tooling}
\begin{itemize}
  \item Problem sets: Python + PyTorch (autograd does the differentiation for you)
  \item You will NOT hand-roll algorithmic differentiation
  \item Fast-computation week (Week 7): C++ exposed to Python via pybind11
  \item Industry pattern: Python for research, C++ for the hot path
\end{itemize}\end{frame}
\begin{frame}{Data, Milestones, Assessment}
\begin{itemize}
  \item Coincall BTC/ETH order books, trades, options chains, funding rates
  \item Milestone 1 (Wk5): Avellaneda-Stoikov engine
  \item Milestone 2 (Wk7): fast pricing \& Greeks library
  \item Milestone 3 (Wk9): full options market-making backtest
  \item Grade: problem sets 30\%, milestones 30\%, final 30\%, participation 10\%
\end{itemize}\end{frame}
\begin{frame}{This Week's To-Do}
\begin{itemize}
  \item Set up Python + PyTorch and a C++17 toolchain with CMake
  \item Build and import the pybind11 'hello world' starter module
  \item Get your Coincall data credentials and sign the data-use agreement
  \item No graded problem set this week
\end{itemize}\end{frame}
\end{document}